\chapter{Conclusiones y Trabajo Futuro}

\section{Conclusiones}

Esta tesis abordó el problema de la reconstrucción de canales EEG faltantes desde la perspectiva del Procesamiento de Señales en Grafos (GSP), con el objetivo de determinar ---bajo condiciones de evaluación rigurosamente imparciales--- si la incorporación de información topológica y temporal produce mejoras reales y significativas frente a los métodos geométricos clásicos. A continuación se presentan las conclusiones principales, organizadas según los objetivos específicos formulados en el Capítulo 1.

\subsection{Validación de la Hipótesis de Trabajo}

La hipótesis planteada fue:
\begin{quote}
\emph{La integración de restricciones de suavidad espacial basadas en topologías de grafos junto con priors de continuidad temporal penalizada permite reconstruir canales faltantes de EEG de manera significativamente más precisa y robusta que los métodos geométricos de extrapolación de superficies estáticas.}
\end{quote}

La evidencia acumulada permite \textbf{confirmar parcialmente y con matices} esta hipótesis. La integración de priors espacio-temporales mejora de forma robusta varias métricas de error frente a implementaciones geométricas propias y frente a MNE-Python en MAE/NRMSE; sin embargo, no domina todas las métricas ni todos los regímenes de señal. La hipótesis queda por tanto apoyada como afirmación condicionada al tipo de métrica, señal, densidad de sensores y presupuesto computacional:

\begin{enumerate}
    \item \textbf{Ventaja observada cuantitativa:} TRSS logró un MAE promedio de \ResTVMAE{} y un SNR de \ResTVSNR~dB en la fase exploratoria, superando a varios baselines geométricos y GSP bajo el protocolo optimizado. En la fase confirmatoria contra \code{MNE interpolate\_bads(method='spline')}, TRSS mantuvo una mejora mediana robusta en MAE de 12.4\% (IC95 7.8--17.4\%) y en NRMSE de 13.9\% (IC95 6.4--19.3\%).

    \item \textbf{La brecha no es puramente paramétrica:} La optimización exhaustiva de baselines con Optuna y la comparación posterior contra MNE-Python muestran que la memoria temporal aporta información complementaria a la interpolación espacial. No obstante, MNE evidencia que una spline esférica bien implementada puede ser muy competitiva cuando la señal es suave, el montaje es denso y la métrica principal no penaliza discontinuidades temporales.

    \item \textbf{Preservación fisiológica con trade-offs:} Los análisis de ERPs y PSD muestran que TRSS puede preservar trayectorias temporales y componentes fisiológicos relevantes bajo pérdidas severas. Sin embargo, la comparación confirmatoria indica que MNE puede conservar mejor la distancia espectral logarítmica agregada en algunos regímenes; por tanto, la preservación fisiológica debe evaluarse de forma multi-métrica y no solo desde error de amplitud.

    \item \textbf{Consistencia multi-dataset no homogénea:} La ventaja de TRSS se observó en datos reales MNE Sample y PhysioNet, además de señales sintéticas rugosas. En señales sintéticas suaves, MNE fue superior. Esta estratificación es valiosa porque delimita el dominio de aplicación: TRSS aporta más cuando existen transitorios, rugosidad espacial o dinámica temporal que una spline instantánea no captura completamente.
\end{enumerate}

\subsection{Cumplimiento de los Objetivos Específicos}

\begin{itemize}
    \item \textbf{Consolidación de Métodos:} Se implementaron y unificaron 18 algoritmos de interpolación y 10 constructores de grafos en un pipeline modular Python, cubriendo familias geométricas, estadísticas, GSP espaciales y temporales.

    \item \textbf{Diseño de Escenarios Clínicos:} Se definieron escenarios experimentales que simulan tanto fallas esporádicas (\textit{random}) como colapsos regionales del amplificador (\textit{nearby}), cubriendo distintos niveles de pérdida. La fase confirmatoria agregó 100 clústeres de escenario dataset $\times$ modo $\times$ severidad para comparar directamente TRSS contra MNE.

    \item \textbf{Optimización Imparcial:} Se integró Optuna con muestreadores TPE y NSGA-II para explorar los espacios paramétricos de todos los métodos ---incluyendo los baselines clásicos--- garantizando que ningún algoritmo fuera evaluado por debajo de su límite de diseño.

    \item \textbf{Evaluación Multi-Métrica:} Se cuantificó el rendimiento con seis métricas complementarias (MAE, RMSE, SNR, DTW, LSD y Coherencia) que capturan dimensiones ortogonales del error, impidiendo que un método pueda ``engañar'' al evaluador mediante sobre-suavizado.

    \item \textbf{Análisis Fisiológico:} Se demostró visual y cuantitativamente el impacto de la reconstrucción en ERPs (N100, P300) y en la distribución de energía por bandas de frecuencia (PSD), proveyendo evidencia directa de relevancia clínica.

    \item \textbf{Trazabilidad:} Se entregó un repositorio con artefactos computacionales versionados (CSV, JSON, figuras, scripts de análisis y tablas LaTeX) que permite auditar los hallazgos principales. La reproducibilidad final requiere congelar dependencias mediante \code{uv.lock}, \code{requirements.txt} o contenedor antes de la entrega institucional.
\end{itemize}

\subsection{Aportes Originales}

Los aportes principales de esta tesis a la comunidad científica son:

\begin{enumerate}
    \item \textbf{Benchmark GSP-EEG con fair tuning y baseline industrial:} Se sometieron métodos novedosos (TRSS, BGSRP) y baselines clásicos (Spherical Spline, RBF) a optimización bayesiana multi-objetivo, y luego se añadió una comparación confirmatoria contra MNE-Python. Esta doble etapa evita evaluar contra baselines débiles y establece un protocolo más exigente para trabajos futuros.

    \item \textbf{Caracterización del cruce de rendimiento:} Se proporcionó evidencia estadística de que la regularización temporal mejora la reconstrucción de amplitud en múltiples escenarios, pero también se identificaron regímenes donde un baseline físico-matemático fuerte como MNE resulta competitivo o superior. El aporte no es una dominancia absoluta, sino una delimitación empírica del cuándo y por qué TRSS aporta.
    
    \item \textbf{Validación a escala trial-level:} Un análisis de 4336 puntos de dato agregados de 736 iteraciones experimentales mediante prueba de Mann-Whitney U confirmó la robustez de los hallazgos: 38 de 55 comparaciones pareadas alcanzaron significancia estadística ($p < 0.001$). Los métodos temporales (temporal\_laplacian, TRSS) y las interpolaciones espaciales simples (IDW, linear) se agrupan en el cuartil superior con MAE mediano $\approx 0.075$-0.081 sin separación significativa, tras excluir el método visibility\_nnk por su comportamiento atípico, demostrando la universalidad del ranking de métodos.

    \item \textbf{Análisis de degradación en casos extremos:} Se observó que bajo regímenes críticos (40\% \textit{nearby}), la regularización temporal puede estabilizar trayectorias fisiológicas que una interpolación espacial instantánea reconstruye de forma menos consistente. Esta afirmación queda limitada a los escenarios y métricas evaluadas.

    \item \textbf{Toolkit de código abierto:} La entrega de un framework modular en Python que unifica carga de datos, construcción de grafos, interpolación, optimización y evaluación, permitiendo a la comunidad replicar y extender los experimentos sobre sus propios datos.
\end{enumerate}

\section{Trabajo Futuro}

Los resultados y limitaciones identificados en esta tesis abren múltiples líneas de investigación futura:

\subsection{Integración BCI en Tiempo Real}

La aplicación más inmediata y de mayor impacto práctico es la integración de TRSS en pipelines BCI \textit{online}. Esto requiere:
\begin{itemize}
    \item \textbf{Ventanas deslizantes:} Adaptar TRSS para operar sobre ventanas temporales solapadas de longitud fija, permitiendo procesamiento causal (solo información pasada y presente).
    \item \textbf{Aceleración GPU:} Paralelizar las operaciones matriciales (producto Laplaciano, gradiente) en GPU mediante frameworks como CuPy o PyTorch, con el objetivo de reducir aún más la latencia medida en el microbenchmark local.
    \item \textbf{Calibración rápida:} Desarrollar heurísticas para la selección rápida de $\alpha$ y $\beta$ basadas en las propiedades espectrales del grafo (e.g., gap espectral), evitando la necesidad de una optimización Optuna completa en tiempo de inferencia.
    \item \textbf{Fallback operativo a MNE:} Diseñar criterios automáticos para escoger entre MNE y TRSS. Por ejemplo, usar MNE cuando la señal es espacialmente suave, el tiempo de cómputo es crítico o la métrica objetivo es espectral; usar TRSS cuando el error de amplitud, la continuidad temporal o fallas contiguas severas son prioritarias.
\end{itemize}

\subsection{Validación Clínica Descendente}

Es necesario evaluar el impacto cuantitativo de TRSS sobre tareas posteriores (\textit{downstream tasks}):
\begin{itemize}
    \item \textbf{Clasificación BCI:} Medir la tasa de acierto de clasificadores de imaginería motora (CSP + LDA o redes convolucionales) cuando los canales faltantes son reconstruidos por TRSS vs. splines vs. descarte.
    \item \textbf{Conectividad funcional:} Evaluar si las matrices de coherencia y correlación parcial computadas sobre señales reconstruidas por TRSS preservan la estructura de red cerebral para diagnóstico de patologías.
    \item \textbf{Poblaciones clínicas:} Validar el framework en datos de pacientes con epilepsia, neonatos y desórdenes de conciencia, donde la dinámica cortical puede diferir sustancialmente de sujetos sanos.
\end{itemize}

\subsection{Modelos Híbridos GSP + Deep Learning}

La combinación de la fundamentación teórica del GSP con la flexibilidad de modelos aprendidos abre posibilidades prometedoras:
\begin{itemize}
    \item \textbf{Graph Neural Networks (GNN):} Entrenar redes neuronales sobre grafos que aprendan simultáneamente la topología óptima y la función de reconstrucción, potencialmente superando la necesidad de hiperparámetros manuales.
    \item \textbf{Autoencoders espaciotemporales:} Utilizar la solución TRSS como inicialización (\textit{warm start}) para autoencoders que refinen la reconstrucción de forma no-lineal.
    \item \textbf{Transfer learning topológico:} Transferir grafos aprendidos en datasets de alta densidad (64 canales) a montajes de baja densidad (22 canales) mediante proyección espectral.
\end{itemize}

\subsection{Extensiones Metodológicas}

\begin{itemize}
    \item \textbf{Degradación gradual:} Extender los protocolos de masking para simular degradación parcial (SNR reducido, no NaN binario) y fallas intermitentes.
    \item \textbf{Regularización adaptativa:} Diseñar esquemas donde $\alpha$ y $\beta$ varíen dinámicamente en el tiempo según la calidad estimada de la señal observada.
    \item \textbf{Criterios espectrales explícitos:} Incorporar LSD, coherencia y potencia de banda directamente en la selección de hiperparámetros de TRSS, dado que la comparación contra MNE mostró que mejorar MAE/NRMSE no garantiza dominar fidelidad espectral.
    \item \textbf{Equivalencia BGSRP:} Cerrar el gap residual entre la implementación Python y la referencia MATLAB/GSPBox para obtener conclusiones definitivas sobre el rendimiento de la familia bandlimited.
    \item \textbf{Grafos temporales dinámicos:} Explorar topologías de grafo que evolucionen con la señal, capturando cambios en la conectividad funcional a lo largo de la sesión.
\end{itemize}

\section{Reflexión Final}

Esta tesis muestra que incorporar conocimiento \textit{a priori} sobre estructura espacial y dinámica temporal puede mejorar la reconstrucción de canales EEG faltantes, especialmente cuando la señal se aleja del régimen suave ideal asumido por las splines esféricas. La comparación contra MNE-Python es central para esta lectura: confirma que TRSS aporta ventajas robustas en error de amplitud, pero también muestra que MNE sigue siendo un baseline físico-matemático fuerte, rápido y difícil de superar en fidelidad espectral agregada.

Más allá del resultado numérico, el aporte metodológico central es la demostración de que el \textit{fair benchmarking} ---incluyendo optimización paramétrica y comparación contra implementaciones de referencia--- es indispensable para establecer conclusiones válidas en la comparación de algoritmos. La contribución de TRSS debe entenderse como complementaria a MNE: una alternativa espacio-temporal prometedora para escenarios dinámicos o adversos, no como sustituto universal de las splines esféricas.
